\section{矩阵的概念}
% ===== 原稿 PDF 第 28 页 =====
% 页首灰底大标题「第二章 矩阵」由章标题自动生成，此处不转录；其右侧并列两个节标题「矩阵的概念／矩阵的运算」

\begin{definition}% 原稿标为「定义 1」
把 $mn$ 个数 $a_{ij}\ (1\le i\le m,\ 1\le j\le n)$ 排成一个矩形阵列，称为一个 $m\times n$ 阶矩阵
\[
A=
\begin{pmatrix}
a_{11} & a_{12} & \cdots & a_{1n}\\
a_{21} & a_{22} & \cdots & a_{2n}\\
\vdots & \vdots & & \vdots\\
a_{m1} & a_{m2} & \cdots & a_{mn}
\end{pmatrix}
=(a_{ij})_{m\times n}.
\]

$a_{ij}$ 称为第 $(i,j)$ 元素。若 $a_{ij}\ (\forall\,i,j)\in\R\ (\C)$，则称为\textbf{实（复）矩阵}。

若 $\forall\,i,j$，$a_{ij}=0$，则称 $A$ 为\textbf{零矩阵}，记为 $O_{m\times n}$。

若 $m=n$，则称 $A$ 为 $n$ 阶\textbf{方阵}。

设 $A=(a_{ij})_{n\times n}$ 为方阵，则称 $a_{11},a_{22},\cdots,a_{nn}$ 连成一线为 $A$ 的\textbf{主对角线}。

若 $\forall\,i\neq j$，$a_{ij}=0$，则称 $A$ 为\textbf{对角阵}，简记为 $\operatorname{diag}\{a_{11},a_{22},\cdots,a_{nn}\}$。
\[
A=
\begin{pmatrix}
a_{11} & 0 & \cdots & 0\\
0 & a_{22} & \cdots & 0\\
\vdots & \vdots & & \vdots\\
0 & 0 & \cdots & a_{nn}
\end{pmatrix}.
\]

进一步，若 $a_{11}=a_{22}=\cdots=a_{nn}=1$，这样的对角阵称为\textbf{单位阵}，记为
\[
I_n=
\begin{pmatrix}
1 & 0 & \cdots & 0\\
0 & 1 & \cdots & 0\\
\vdots & \vdots & & \vdots\\
0 & 0 & \cdots & 1
\end{pmatrix}.
\]

若 $a_{ij}=0\ (\forall\,i>j)$，则称 $A$ 为\textbf{上三角阵}，若 $a_{ij}=0\ (\forall\,i<j)$，则称为\textbf{下三角阵}。
\end{definition}

\begin{definition}% 原稿标为「定义 2」
设 $A=(a_{ij})_{m\times n}$，$B=(b_{ij})_{s\times t}$。
\[
A=B\overset{\mathrm{def}}{\Longleftrightarrow}m=s,\quad n=t,\quad a_{ij}=b_{ij}
\quad\left(\substack{\forall\,1\le i\le m\\[1pt] \forall\,1\le j\le n}\right).
\]
\end{definition}

\begin{example}% 原稿标为「例」
\[
(1)\ \begin{pmatrix}
0 & 0\\
0 & 0
\end{pmatrix}_{2\times2}
\neq
\begin{pmatrix}
0 & 0 & 0\\
0 & 0 & 0
\end{pmatrix}_{2\times3}
\qquad\qquad
(2)\ \begin{pmatrix}
1 & 0\\
0 & 0
\end{pmatrix}
\neq
\begin{pmatrix}
1 & 1\\
0 & 0
\end{pmatrix}
\]
\end{example}

\begin{definition}% 原稿标为「定义 3」
设 $A=(a_{ij})_{m\times n}$，则 $(a_{i1},a_{i2},\cdots,a_{in})$ 称为 $A$ 的第 $i$ 个\textbf{行向量}；
\[
\begin{pmatrix}
a_{1j}\\
a_{2j}\\
\vdots\\
a_{mj}
\end{pmatrix}
\]
称为 $A$ 的第 $j$ 个\textbf{列向量}。

称 $1\times n$ 矩阵 $(a_1,a_2,\cdots,a_n)$ 称为 $n$ 维行向量。

称 $n\times 1$ 矩阵 $\begin{pmatrix}a_1\\a_2\\\vdots\\a_n\end{pmatrix}$ 称为 $n$ 维列向量。
\end{definition}

记 $M_n(\R)$ 为 $n$ 阶实方阵全体构成的集合（$M_n(\C)$ 为 $n$ 阶复方阵全体构成的集合）。
% 原稿在 M_n(R) 的 R 下注「(C)」、在「实」下注「(复)」，即同时给出复矩阵的情形

映射（函数）：
\[
\begin{array}{r@{\ }c@{\ }l}
\det\colon & M_n(\R) & \longrightarrow\ \R\\[2pt]
& A & \longmapsto\ |A|=\det(A)
\end{array}
\]

% ===== 原稿 PDF 第 29 页 =====

\noindent\textbf{问题}\quad (1) $n$ 阶行列式的值在多大程度上反映 $n$ 阶方阵的性质？

(2) 映射 $\det$ 具有怎样的性质？

\section{矩阵的运算}% 编者所加：原稿无独立标题

\noindent\textbf{矩阵的加减法}\quad 设 $A=(a_{ij})_{m\times n}$，$B=(b_{ij})_{m\times n}$。
\[
\begin{aligned}
A+B&\overset{\mathrm{def}}{=}(a_{ij}+b_{ij})_{m\times n},\\
A-B&\overset{\mathrm{def}}{=}(a_{ij}-b_{ij})_{m\times n}.
\end{aligned}
\]

\begin{example}% 原稿标为「例」
\[
(1)\ \begin{pmatrix}
1 & 2 & 0\\
-1 & 3 & 1
\end{pmatrix}
+
\begin{pmatrix}
1 & -2 & 1\\
-1 & -3 & 0
\end{pmatrix}
=
\begin{pmatrix}
2 & 0 & 1\\
-2 & 0 & 1
\end{pmatrix}
\]
\[
(2)\quad A_{m\times n}\pm O_{m\times n}=A_{m\times n}
\]
\end{example}

\begin{definition}[加法性质]% 原稿标为「定义 4 加法性质」
\[
\begin{array}{r@{\quad}l@{\qquad}l}
(1) & \text{交换律：} & A+B=B+A,\\
(2) & \text{结合律：} & (A+B)+C=A+(B+C),\\
(3) & \text{零矩阵：} & A+O=A,\\
(4) & \text{负矩阵：} & A+(-A)=O,\\
(5) & & A-B=A+(-B).
\end{array}
\]
\end{definition}

\begin{definition}[数乘]% 原稿标为「定义 5」
设 $A=(a_{ij})_{m\times n}$，$c$ 为常数。% 原稿作「C 为常数」，下文公式中均记作 c

$c$ 与 $A$ 的数乘为（scalar product）
\[
c\cdot A\overset{\mathrm{def}}{=}(c\cdot a_{ij})_{m\times n}.
\]
负矩阵
\[
-A\overset{\mathrm{def}}{=}(-1)\cdot A=(-a_{ij})_{m\times n}.
\]
\end{definition}

\begin{example}% 原稿标为「例」
\[
2\cdot
\begin{pmatrix}
1 & 2\\
3 & 4
\end{pmatrix}
=
\begin{pmatrix}
2 & 4\\
6 & 8
\end{pmatrix}.
\]
\end{example}

\noindent\textbf{数乘的性质}
\[
\begin{array}{r@{\quad}l@{\qquad}l}
(1) & \text{数的分配律：} & (c+d)A=cA+dA,\\
(2) & \text{矩阵的分配律：} & c\cdot(A+B)=c\cdot A+c\cdot B,\\
(3) & \text{数乘的结合律：} & (cd)\cdot A=c\cdot(d\cdot A),\\
(4) & \text{数乘单位元：} & 1\cdot A=A,\\
(5) & \text{数乘零元：} & 0\cdot A_{m\times n}=O_{m\times n}.
\end{array}
\]

\begin{definition}[矩阵乘法]% 原稿标为「定义 6 矩阵乘法」
设 $A=(a_{ij})_{m\times k}$，$B=(b_{ij})_{k\times n}$，定义 $A$ 与 $B$ 的乘积 $C$ 是 $m\times n$ 矩阵，其元素
\[
c_{ij}=a_{i1}b_{1j}+a_{i2}b_{2j}+\cdots+a_{ik}b_{kj}=\sum_{r=1}^{k}a_{ir}b_{rj}.
\]
% 原稿求和下标手写略似「n=1」，按通项 a_{ir}b_{rj} 记为 ∑_{r=1}^{k}
\end{definition}

\begin{remark}% 原稿标为「注」
(1) 只有 $A$ 的列数 $=$ $B$ 的行数，$AB$ 才有意义。
\end{remark}
